\usepackage[spanish,es-tabla]{babel}
\usepackage[top=2cm,bottom=2cm,right=2cm,left=2cm]{geometry}
\usepackage{amsmath}
\usepackage{amssymb}
\usepackage{amsthm}
\usepackage{empheq}
\usepackage{mdframed}
\usepackage{booktabs}
\usepackage{lipsum}
\usepackage{graphicx}
\usepackage{color}
\usepackage{psfrag}
\usepackage{pgfplots}
\usepackage{bm}
\usepackage{multicol}
\usepackage{wrapfig}
\usepackage{float}
\usepackage{wrapfig}
\usepackage{multirow}
\usepackage{biblatex}
\usepackage{csquotes}
\usepackage{parskip}
\usepackage{cancel}
\usepackage{listing}
\usepackage{tikz}
\usepackage{verbatim}
\usepackage{subcaption}
\usepackage{enumitem}
\usepackage{xcolor}
\usepackage{xfrac}
\usepackage[load-configurations = abbreviations]{siunitx}
\usepackage{fix-cm}
\usepackage{yhmath}
\usepackage[bookmarks=true,colorlinks=true,linkcolor=black,citecolor=black,menucolor=black,urlcolor=black,hidelinks,breaklinks=true]{hyperref}
\usepackage{mathrsfs}
\usepackage{tcolorbox}
\usepackage{amsfonts}
\usepackage[bb=ams]{mathalpha}
\usepackage{mathrsfs}
\usepackage{xstring}

\graphicspath{{fig/}}
\addbibresource{biblio.bib}
\usepgfplotslibrary{colorbrewer}
\pgfplotsset{width=8cm,compat=1.9}
\definecolor{negro}{RGB}{0,0,0}
\definecolor{azul}{RGB}{0,0,255}
\definecolor{verde}{RGB}{0,255,0}
\definecolor{cian}{RGB}{0,255,255}
\definecolor{rojo}{RGB}{255,0,0}
\definecolor{magenta}{RGB}{255,0,255}
\definecolor{amarillo}{RGB}{255,255,0}
\definecolor{blanco}{RGB}{255,255,255}

\definecolor{limegreen}{cmyk}{0.75,0,1,0}
\definecolor{orange}{cmyk}{0,0.8,0.95,0}
\definecolor{purple}{cmyk}{0.75,1,0,0}

\newcommand{\uvec}[1]{~ \boldsymbol{\hat{\textbf{#1}}}}

\newsavebox\foobox
\newlength{\foodim}
\newcommand{\slantbox}[2][0]{\mbox{%
        \sbox{\foobox}{#2}%
        \foodim=#1\wd\foobox
        \hskip \wd\foobox
        \hskip -0.5\foodim
        \pdfsave
        \pdfsetmatrix{1 0 #1 1}%
        \llap{\usebox{\foobox}}%
        \pdfrestore
        \hskip 0.5\foodim
}}
\def\Laplacesymb{\slantbox[-.45]{$\mathscr{L}$}}
\def\inverseLaplacesymb{\Laplacesymb^{-1}}
\def\Fouriersymb{\slantbox[-.45]{$\mathscr{F}$}}
\def\Fouriersinsymb{\Fouriersymb_{\textup{s}}}
\def\Fouriercossymb{\Fouriersymb_{\textup{c}}}
\def\inverseFouriersymb{\Fouriersymb^{-1}}

\newcommand{\Laplace}[2]{
    \IfEqCase{#1}{
        {}{\Laplacesymb\left\lbrace#2\right\rbrace}
        {i}{\inverseLaplacesymb\left\lbrace#2\right\rbrace}
    }[\PackageError{tree}{Undefined option to tree: #1}{}]
}
\newcommand{\Fourier}[2]{
    \IfEqCase{#1}{
        {}{\Fouriersymb\left\lbrace#2\right\rbrace}
        {i}{\inverseFouriersymb\left\lbrace#2\right\rbrace}
        {sin}{\Fouriersinsymb\left\lbrace#2\right\rbrace}
        {cos}{\Fouriercossymb\left\lbrace#2\right\rbrace}
    }[\PackageError{tree}{Undefined option to tree: #1}{}]
}

\sisetup{
  detect-family,
  per-mode = symbol,
  separate-uncertainty = true,
  inter-unit-product = \cdot,
  %output-decimal-marker = {,},
  %exponent-product = \cdot
}

\begin{document}

\onecolumn\begin{@twocolumntrue}
    \begin{minipage}{0.3\textwidth}{\includegraphics[scale=0.24]{Escudo_UD.pdf}} % Cambiar escudo
    \end{minipage}
    \vspace{10pt}
    \begin{minipage}{0.677\textwidth}
        \begin{center}
            \vspace{12mm}
        
            \Large{\textbf{<<Título>>}}
            \vspace{3mm}
            
            \large{\textbf{<<Autores>>}}
            \vspace{2mm}
    
            \large{\textit{<<Carrera>>, <<Universidad>>}}
            \vspace{1mm}
            
            <<Fecha>> % Formato (mes) de (año)
        \end{center}
\vspace{5pt}
\end{minipage}

\centerline{\rule{0.95\textwidth}{0.4pt}} % Línea negra
\end{@twocolumntrue}

\begin{multicols}{2}

\section{Periodo de una función}

\begin{enumerate}
    \item Encontrar el periodo de las siguientes funciones:
    \begin{enumerate}
        \item $f(t) = cos \left(\pi t\right)$ 
        \item $f(t) = \sin \left(t\right) + \sin \left(\frac{t}{3}\right) + \sin \left(\frac{t}{5}\right)$
    \end{enumerate} 
\end{enumerate}

\subsubsection*{\textbf{Solución:}}


Para la solución de este ejercicio, se tiene en cuenta la forma general para hallar el periodo de una función trigonométrica, por tanto, para toda función trigonométrica con periodo natural $T_n$, se tiene su periodo $T_f$ de la siguiente manera:
\begin{equation}
    T_f = \frac{T_n}{B}
\label{eq1}
\end{equation}
Dónde $B$ yace de la forma general de una función trigonométrica, la cuál es:
\begin{equation*}
    f(t) = A\sin\left(Bx + C\right) + D
\end{equation*}
Forma usada para cualquier función trigonométrica no necesariamente para la función seno, de esta manera, solo analizando la forma de la función y además conociendo el periodo natural de la función que estemos utilizando, podremos conocer el periodo de dicha función.

Los periodos naturales de las diferentes funciones trigonométricas son:

\begin{itemize}
    \item Para \textit{seno, coseno, secante} y \textit{cosecante} es $2\pi$.
    \item Para \textit{tangente} y \textit{cotangente} es $\pi$.
\end{itemize}

Para una superposición de funciones, se debe buscar el mínimo común múltiplo (m.c.m.) de esta manera se busca en qué instante todas las funciones pertenecientes a la superposición están en fase, o lo que sería igual a que comparten el mismo periodo.

Teniendo todo esto en cuenta y aplicándolo, resulta:

$a) ~ f(t) = \cos \left(\pi t\right)$

Tomando (\ref{eq1}) y teniendo $B = \pi$ y $T_n = 2\pi$, el periodo $T_f$ de la función es:
\begin{equation}
    T_f = \frac{2 \cancel{\pi}}{\cancel{\pi}} \Rightarrow \textcolor{rojo}{T_f = 2}
    \label{res1}
\end{equation}
Obteniendo así que el periodo de la función es \textcolor{rojo}{$T_f = 2$}.

$b) ~ f(t) = \sin \left(t\right) + \sin \left(\frac{t}{3}\right) + \sin \left(\frac{t}{5}\right)$

Tomando (\ref{eq1}) nuevamente y teniendo $B_1 = 1$, $B_2 = \left(\frac{1}{3}\right)$, $B_3 = \left(\frac{1}{5}\right)$ y $T_n = 2\pi$ para todas las funciones superpuestas, definimos $T_f$ para cada una y posteriormente hallamos el m.c.m. entre ellas, así de esta manera obtener el periodo general de la función, por tanto:
\begin{align}
    T_{f1} & = \frac{2\pi}{1} \Rightarrow T_{f1} = 2\pi \\
    T_{f2} & = \frac{2\pi}{ \frac{1}{3} } \Rightarrow T_{f2} = 6\pi \\
    T_{f3} & = \frac{2\pi}{ \frac{1}{5} } \Rightarrow T_{f3} = 10\pi
\end{align}
Usando los resultados obtenidos en (3), (4) y (5)
\begin{center}

$\begin{array}{c c c | c}
    2\pi & 6\pi & 10\pi & 2\\
    1\pi & 3\pi & 5\pi & 3\\
    \hphantom{1\pi} & 1\pi & 5\pi & 5\\
    \hphantom{1\pi} & \hphantom{1\pi} & 1\pi & \hphantom{1}\\ \hline
    \hphantom{1\pi} & \hphantom{1\pi} & \hphantom{1\pi} & 30
\end{array}$

\end{center}
De esta manera se concluye qué:
\begin{equation}
    \textcolor{rojo}{T_f = 30\pi}
\label{res2}
\end{equation}

\begin{enumerate}[resume]
    \item Si $f(t)$ es una función  periódica de $t$ con periodo $T$, demostrar que $f(at)$, para $a \neq 0 $, es una función periódica de $t$ con periodo $\frac{T}{a}$.
\end{enumerate}

\subsubsection*{\textbf{Solución:}}

A la hora de demostrar que una función $f(t)$ posee un periodo $T$ y lo que implicaría que tiene la forma $f(t+T)$

\section{Sección 1} 

\begin{equation}
    Ecuaciones
\label{ecuacion}
\end{equation}

\begin{equation*} % Para ecuaciones sin número y que el igual se mantenga en la misma sangría
\begin{split}
    texto1 & = texto---2 \\
    texto---2 & = texto------3 \\
    texto------3 & = texto4
\end{split}
\end{equation*}

\section{Sección 2}

<<Texto>>

\end{multicols}

\subsection{Subsección 1}

<<Variación entre dos y una columna>> -------------------------------------------------------------------

\begin{multicols}{2}

\subsection{Subsección 2}

<<Texto>>

\subsubsection{Subsubsección 1}

<<Texto>>

\begin{figure}[H]
    \centering
    \includegraphics[scale=0.32]{pollito.png}
    \caption{Inserción de una figura}
    \label{figura}
\end{figure}

\section{Conclusiones}

<<Conclusión>>

\printbibliography

\end{multicols}

\section{Anexos}

<<Anexo>>

\end{document}